\chapter{Discussion}
% A particularly important chapter is the discussion of the approach and the results. The aim here is to critically investigate the quality of your solution: what are the potentials, and more importantly, what are the limitations? How do your findings relate to previous work? To what extent will your results generalize? In essence: what is the impact of your work?



\section{Main Findings}
% Shortly summarize your main findings by answering each research question in turn.
\textbf{\ac{RQ}1: Which forecasting approach provides the best predictive performance for short-term water level forecasting given the available data?}\\
\ac{RNN} achieves the best \ac{CV} \ac{RMSE} with a value of $12.02$ cm, slightly outperforming other models like XGBoost or \ac{MLP}. However, XGBoost has the lowest standard deviation of \ac{RMSE} across the \ac{CV} folds. On the sealed test set \ac{RNN} does not achieve the lowest \ac{RMSE}, but the \ac{MLP} model does, with an \ac{RMSE} of $12.11$ cm compared to $12.85$ cm for Ridge and $12.96$ cm for \ac{RNN}. Therefore, \ac{MLP} is selected as the final model, since the final model selection is based on test set performance, while \ac{CV} is used to select the feature subsets and hyperparameters for each model family. Compared to the persistence baseline, with an \ac{RMSE} of $16.00$ cm, this corresponds to an improvement of about $24.3\%$. The better performance of the \ac{MLP} indicates that this model generalized better than the other models to this particular test period, which includes previously unseen extreme water levels. However, based on a single test period it cannot be concluded that \ac{MLP} generally performs better than other models. Therefore, based on the predefined model-selection criterion, \ac{MLP} provides the best predictive performance for the investigated test period, although \ac{RNN} achieves the best performance in \ac{CV}.
\\ \\
\textbf{\ac{RQ}2: How does incorporating data from upstream stations and weather information affect forecast performance?}\\
The target station alone feature subset without weather information achieves the worst \ac{CV} \ac{RMSE} across the models evaluated on different feature subsets. The addition of upstream stations and weather information improves the \ac{CV} \ac{RMSE} across these models. For \ac{MLP} the \ac{CV} \ac{RMSE} improves by about $23.1\%$ from $16.44$ cm (just the target station) to $12.64$ cm (all features). When just looking at either upstream stations or weather information, the upstream stations lead to a larger improvement in \ac{CV} \ac{RMSE} than the weather information alone, with an improvement of $12.3\%$ compared to $3.2\%$ for \ac{MLP}. The combination of both upstream stations and weather information leads to the best performance for \ac{MLP} and Ridge, indicating that both sources provide additional predictive information. However, for \ac{RNN} the best performance is achieved when only using the current water levels from all stations, indicating that more features do not always lead to better performance. The improved performance from including upstream stations indicates that their water level measurements contain relevant information about future water levels at the target station, likely because changes in water level propagate downstream. Furthermore, the additional improvement from weather information for \ac{MLP} and Ridge indicates that it provides information that is not captured by the upstream stations alone. Therefore, it is beneficial to include both upstream stations and weather information for short-term water level forecasting with these models.
\\ \\
\textbf{\ac{RQ}3: How does forecast performance vary with lead time and across different water level regimes?}

With increasing lead time, the forecast error of \ac{MLP} and Ridge generally increases. The persistence model has an advantage over \ac{MLP} only at H+1, while Ridge performs better than persistence across the entire forecast horizon. Ridge performs better than \ac{MLP} for the first two hours. From H+3 \ac{MLP} has a lower \ac{RMSE} per forecast horizon compared to Ridge, indicating that it performs better at longer lead times. When looking at the different water level regimes, both models perform well for the first three quartiles, but the performance deteriorates for the highest water levels. Both models especially perform poorly when looking at forecasted values where the ground truth is at or above the alarm threshold of 545 cm. With an \ac{ME} of $-23.51$ cm and an \ac{RMSE} of $47.33$ cm for \ac{MLP}, compared to $-14.09$ cm and $47.92$ cm for Ridge, both models systematically underestimate the extreme water levels on average. However, both models are able to capture the overall trend of the September 2024 flood at H+1, although this does not mean that they perform equally well at longer lead times. This highlights a limitation of the models in capturing rare but critical events across the full forecast horizon.
\section{Comparison with Previous Work}
% An important part of this chapter is the comparison with existing solutions and approaches found in the literature. How do the findings of the thesis relate to existing approaches? Are the results plausible? Is the proposed solution better than existing approaches?

The finding that no single model performs best across all evaluations is consistent with previous work. \citeauthor{asifReviewIntercomparisonMachine2025} show that model performance changes with the forecast horizon, with Random Forest performing well between 12 and 48 hours \cite{asifReviewIntercomparisonMachine2025}. Furthermore, \citeauthor{belyakovaForecastingWaterLevels2022} find that XGBoost performs best for the Pshish River, while \ac{MLP} performs best for the Mzymta River \cite{belyakovaForecastingWaterLevels2022}. In this thesis, \ac{RNN} achieves the lowest \ac{CV} \ac{RMSE}, while \ac{MLP} performs best on the sealed test set. The \ac{CV} \ac{RMSE} of XGBoost, \ac{RNN}, and \ac{MLP} differs by at most $0.62$ cm, indicating that their average performance is relatively close. \ac{MLP} also performs best overall in the Muda River study, although that study uses daily data \cite{zakariaExploringMachineLearning2023}. Thus, the selection of \ac{MLP} is plausible, but does not mean that it generally performs better than the other approaches.

When comparing the absolute errors, the forecast horizon and evaluation data have to be considered. The Parma River, tidal river, and Krasnodar Krai studies evaluate shorter horizons, while the operational system of \citeauthor{nevoFloodForecastingMachine2022} has maximum lead times of eight to 48 hours and therefore covers a comparable forecast horizon \cite{dazziFloodStageForecasting2021,guoPredictionRiverStage2021,belyakovaForecastingWaterLevels2022,nevoFloodForecastingMachine2022}. However, the $12.11$ cm test-set \ac{RMSE} of \ac{MLP} cannot be directly compared with their results. For example, \citeauthor{dazziFloodStageForecasting2021} evaluate extracted flood events, while this thesis evaluates eligible forecasts from a continuous record across the full 24-hour horizon. Furthermore, the rivers and performance measures differ. Therefore, the absolute errors do not show that the proposed model is better or worse than previous models.

The positive effect of upstream stations and weather data is consistent with \citeauthor{nevoFloodForecastingMachine2022}. Their models also use past and upstream stage measurements, and adding precipitation improves the \ac{LSTM} at $89\%$ of the gauges \cite{nevoFloodForecastingMachine2022}. This is similar to this thesis, where upstream stations reduce the \ac{MLP} \ac{CV} \ac{RMSE} by $12.3\%$ and weather data reduce it by $3.2\%$, each compared to using only the target station without weather data. Adding both leads to an improvement of $23.1\%$. However, additional data do not always help. The Muda River study reports no significant improvement from adding meteorological data \cite{zakariaExploringMachineLearning2023}. Similarly, \citeauthor{sandersDataDrivenFloodAlert2022} find no overall improvement from rain gauges outside the watershed, while these gauges delay the predicted peak \cite{sandersDataDrivenFloodAlert2022}. In this thesis, \ac{RNN} performs best when only using the current water levels from all stations. Therefore, the benefit of additional data depends on both the catchment and the model.

The strong performance of persistence at short lead times is also consistent with previous work. On the Pshish test set, XGBoost performs worse than persistence at one hour, but considerably better at 15 hours \cite{belyakovaForecastingWaterLevels2022}. In this thesis, persistence performs better than \ac{MLP} only at H+1. \citeauthor{nevoFloodForecastingMachine2022} discuss the slow response of large rivers and the strong performance of persistence compared to predicting the mean \cite{nevoFloodForecastingMachine2022}. This is also a plausible explanation for the Danube, where the current water level provides a useful forecast for the next hour. However, Ridge performs better than persistence across the entire forecast horizon. Thus, persistence is a useful baseline, but its advantage at short lead times also depends on the model.

The poor performance at extreme water levels is also consistent with previous work. \citeauthor{dazziFloodStageForecasting2021} report increasing flood-peak errors at longer lead times \cite{dazziFloodStageForecasting2021}. \citeauthor{guoPredictionRiverStage2021} find underestimation of high water levels in training results at two stations, but also peak overestimation in validation \cite{guoPredictionRiverStage2021}. These findings show difficulties with high water levels, although the errors do not always have the same direction. In this thesis, \ac{MLP} and Ridge underestimate water levels at or above the alarm threshold on average across the full forecast horizon. However, both models capture the overall trend of the September 2024 flood at H+1, although the flood exceeds the training range and \ac{MLP} overestimates the highest peak. Therefore, the poor average performance at extreme water levels does not mean that every forecast underestimates them.

The evaluation by water level regime shows that the overall \ac{MLP} \ac{RMSE} of $12.11$ cm hides an increase to $47.33$ cm at or above the alarm threshold. Similarly, \citeauthor{sandersDataDrivenFloodAlert2022} report almost identical overall errors for two input scenarios, although the peak prediction for one event changes from five to 85 minutes late \cite{sandersDataDrivenFloodAlert2022}. This shows that good overall performance does not mean good performance during critical events. Therefore, the evaluation of individual flood events and high water levels provides relevant information in addition to the overall error. For the intended use in FloodAlert, the results support displaying the forecast as additional information instead of using it as the sole basis for automated warnings.

Overall, the findings agree with previous work in showing that forecast quality depends on the river, input data, model, and forecast horizon. However, hybrid approaches, which perform well in the review by \citeauthor{asifReviewIntercomparisonMachine2025}, are not evaluated in this thesis \cite{asifReviewIntercomparisonMachine2025}. Therefore, the results support the selection of \ac{MLP} among the evaluated models for Korneuburg and the investigated period. The evaluation procedure can also be applied to other stations, but the model and feature subsets would have to be evaluated again.



\section{Implications}
\todo{change this to match the latest findings}
% Also, you should assess impacts on society and/or the company (whatever applies) – positive and negative ones. This could include benefits such as organizational improvements (cost reductions), increased benefits (for new business models), or threats such as loss of jobs, increase of faked information, biases in automated information handling, loss of privacy, legal/ethical issues, etc.
% Sending out future alters should be done with caution.
The main practical impact on society and the company is the ability to provide additional information about the future water levels. However, the accuracy of the forecasts is especially weak for extreme events, such as the September 2024 flood, where the model systematically underpredicts the water levels. Therefore, caution should be exercised when deciding on actions based on these forecasts, like sending out early warnings, or promoting the feature of early flood warnings, based on forecasts. Especially for water levels above the alarm threshold, the model shows a strong systematic underprediction, as also illustrated in \autoref{fig:mlp-worst-forecast-window}. The performance of the current model is therefore not sufficient to use the forecasts as the sole basis for flood warnings or other critical decisions. Instead, the forecasts could provide additional information alongside existing measurements and warning mechanisms.

\section{Limitations}
% Weaknesses in the methodology should be listed (e.g., sample sizes, data quality, limited choice of models, non-systematic hyperparameter tuning, incomplete implementation of the deployment, performance limitations etc.) and the impact on the solution be discussed. Also, possible improvements should be suggested. It is mandatory to comment on possible bias in data and/or models (e.g., gender, ethnicity …).

% Single target station only 

% Rare extreme events
There are several limitations that need to be addressed: Firstly the limited generalizability of the model, as it has been only trained on a single target station. Furthermore the struggle with accurately predicting rare extreme events.

Only using a single target station for training results in limited generalizability, as the model has not been tested on other stations or different rivers with varying hydrological characteristics. This implies that the performance of the model may significantly vary when applied to new locations. A possible improvement would be to train and evaluate the model on multiple stations and rivers, also with different hydrological characteristics (e.g. different sizes) to get an overall better understanding of its generalizability and robustness.

Additionally water levels above the alarm threshold of 545 cm are particularly challenging for the model to predict accurately, since there are relatively few instances, only about $0.76\%$ of the directly observed hourly water levels in the dataset. This scarcity of events makes it difficult for the model to learn the underlying patterns and may contribute to systematic underprediction for extreme water levels. Both \ac{MLP} and Ridge underestimate water levels at or above the alarm threshold on average across the full forecast horizon. However, this does not mean that every forecast underestimates them, as \ac{MLP} overestimates the highest peak of the September 2024 flood at H+1. A possible mitigation for this could be to train a binary classifier to predict whether there will be an extreme event, and then use a specialized model to predict the exact water level for those events.

Another limitation is that the models only use historical and current weather observations and do not incorporate weather forecasts. Especially precipitation forecasts could provide valuable information for predicting future water levels more accurately. This could help models to anticipate future changes in water levels and may improve the overall forecast performance. However, this would need to be evaluated, since weather forecasts also contain errors.

Potential demographic biases, such as gender or ethnicity, are not relevant for the water level sensor data used in this thesis. However, the data is biased towards normal water level conditions, since extreme water levels are rare. Since only a single target station on the Danube is evaluated, the results are biased towards the hydrological conditions at this location.

\section{Further Research}
% Ideas: Compare Transformer models e.g. Temporal Fusion Transformer (TFT). Add not only the upstream rivers from the same river but also upstream rivers that merge into the target river. Compare the performance of different models on different rivers.
% Use different models for different regimes.
There are many approaches for further research in this area. One direction could be to try to improve the forecast performance for extreme water levels. This could be done by developing specialized models for extreme events or different regimes, such as low, normal, and high water levels. Furthermore, the inclusion of the decades of historical data (in daily resolution) could help to provide additional extreme event examples. However, the daily resolution would need to be considered for the hourly forecasting task. Future research could also investigate the use of more complex models such as Transformer-based architectures like \ac{TFT}.

Another potential for future research is to explore the generalizability of the model on other measurement stations, by applying the same forecasting pipeline to additional stations. This would help to assess whether \ac{MLP} maintains its performance across different locations and hydrological conditions, whether other models perform better at these locations and whether the addition of upstream and weather data improves the performance consistently. Additionally, including weather forecast data, not just the historical observations, could further enhance the predictive capabilities of the model.
